% training.tex -- Section 11: Training Pipeline

\section{Training Pipeline}
\label{sec:training}

Kant includes a training pipeline that uses game-theoretic environments to
improve language model alignment through reinforcement learning. This section
describes the GRPO and DPO training methods, self-play training, curriculum
learning strategy, and reward shaping approach.

The training pipeline tests whether game-theoretic experience transfers to
broader alignment. We use both online (GRPO) and offline (DPO) methods,
with the Kant environment as the reward oracle.

\subsection{GRPO on Game Environments}
\label{sec:grpo}

Group Relative Policy Optimization (GRPO)~\citep{shao2024deepseekmath}
generates $G$ completions per prompt and optimizes the policy by comparing
relative rewards within each group. We implement GRPO via TRL with the
following configuration:

\begin{itemize}[nosep]
    \item \textbf{Group size}: $G = 8$ completions per game state
    \item \textbf{Max completion length}: $16$ tokens (moves are $1$--$2$ words)
    \item \textbf{Learning rate}: $5 \times 10^{-6}$
    \item \textbf{Batch size}: $4$ with gradient accumulation $\times 4$
    \item \textbf{Training steps}: $500$ (across $1{,}000$ game prompts)
    \item \textbf{Base models}: 8 open-weight models from 1B to 27B parameters
          (see Section~\ref{sec:safety-setup})
\end{itemize}

\noindent Each prompt presents a game state---name, description, available
moves, opponent strategy, round number, history, and cumulative score.
The LLM completion is parsed as a move and submitted to the Kant environment;
the returned payoff becomes the GRPO reward signal. Invalid or unparseable
moves receive reward $-2.0$; environment errors default to $0.0$.

\subsection{DPO via Trajectory Preference Pairs}
\label{sec:dpo}

Direct Preference Optimization (DPO)~\citep{rafailov2023direct} trains on
preference pairs constructed from game trajectories. For each game episode,
we generate trajectory pairs where:

\begin{itemize}[nosep]
    \item \textbf{Preferred} ($y_w$): cooperative, fair, strategically sound
          moves (e.g., cooperate in PD against TFT, fair offers in Ultimatum)
    \item \textbf{Rejected} ($y_l$): exploitative, unfair, or naive moves
          (e.g., always defect, lowball offers, free-riding)
\end{itemize}

\noindent The DPO loss $\mathcal{L}_{\text{DPO}} = -\log\sigma\bigl(
\beta \log \frac{\pi_\theta(y_w|x)}{\pi_{\text{ref}}(y_w|x)} -
\beta \log \frac{\pi_\theta(y_l|x)}{\pi_{\text{ref}}(y_l|x)}\bigr)$
directly optimizes the policy to prefer aligned trajectories without an
explicit reward model. The preference ordering
``cooperative $\succ$ exploitative'' encodes the alignment objective.

\subsection{Self-Play Training}
\label{sec:self-play}

To move beyond fixed opponent strategies, Kant supports self-play training
where the agent plays against frozen copies of itself or a pool of past
checkpoints.

\paragraph{FrozenOpponent.}
A \texttt{FrozenOpponent} wraps a generation function for inference without
gradients. Opponents can be instantiated from the current model
(\texttt{from\_model}), a saved checkpoint (\texttt{from\_checkpoint}), or
an external API (\texttt{from\_api}) for cross-model self-play.

\paragraph{OpponentPool.}
The \texttt{OpponentPool} maintains a collection of past model checkpoints
as diverse opponents. New checkpoints are added at a configurable update
interval; when the pool exceeds its maximum size, oldest entries are evicted
(FIFO). During training, opponents are sampled uniformly at random from the
pool, preventing overfitting to any single opponent style.

\paragraph{Cross-model self-play.}
The API-based opponent interface enables self-play against models served via
external inference endpoints, supporting heterogeneous training where agents
learn from opponents with fundamentally different architectures and scales.

\subsection{Curriculum Learning}
\label{sec:curriculum}

Training proceeds through a curriculum that progressively expands the game
domain:

\begin{enumerate}[nosep]
    \item \textbf{Phase~I}: Classical dilemmas (PD, Stag Hunt, Hawk-Dove)
    \item \textbf{Phase~II}: Extended matrix and sequential games
    \item \textbf{Phase~III}: N-player and coalition games
    \item \textbf{Phase~IV}: Governance-enabled games and dynamic creation
\end{enumerate}

Each phase doubles the game pool, with transition triggered by convergence
of the composite strategic reasoning metric $M_S > 0.6$ on the current pool.

\subsection{Reward Shaping}
\label{sec:reward-shaping}

Raw game payoffs are transformed to balance individual performance with
cooperative and fairness objectives:
\[
r_{\text{shaped}} = \alpha \cdot u_i + \beta \cdot (u_i + u_{-i})
  - \gamma \cdot |u_i - u_{-i}| - \delta \cdot \mathbb{1}[\text{exploit}]
\]
where $u_i$ is the player's payoff, $u_{-i}$ the opponent's, and the
exploit indicator fires when the agent achieves $T$ (temptation payoff)
against a cooperative opponent. Default weights:
$\alpha = 0.4$, $\beta = 0.3$, $\gamma = 0.2$, $\delta = 0.1$.
This creates a gradient toward Pareto-efficient, fair outcomes while
maintaining strategic competence.
